\section*{Original search, assignment and branching-program definitions (pp.17:2,17:4--6)}
Consider a search problem, defined by a relation $S\subseteq I\times O$ for some finite sets $I$ and $O$. On input $x\in I$ the search problem is to find some output in $S(x):=\{o\in O\mid(x,o)\in S\}$.
\begin{result}{Definition 1}
A unsatisfied search problem $\Search_\varphi$ for an unsatisfiable CNF formula $\varphi:=\bigwedge_{i\in I}C_i$ on $n$ variables is defined as follows:
\begin{itemize}\item input: an $n$-variable assignment $z\in\{0,1\}^n$;
\item output: an element $i\in I$ such that clause $C_i$ of $\varphi$ is falsified by $z$.\end{itemize}
\end{result}
Let $X$ be a set of boolean variables. For a variable $x\in X$ we denote $x^1:=x$ and $x^0:=\neg x$.
We say that $\alpha:X\to\{0,1,*,?\}$ is a generalized partial assignment and $\alpha$ assigns or touches $x\in X$ iff $\alpha(x)\in\{0,1,?\}$. And an assignment $\gamma$ is an instance of $\alpha$ iff:
\begin{itemize}\item $\alpha(x)\in\{0,1,*\}$ implies $\gamma(x)=\alpha(x)$;
\item $\alpha(x)=?$ implies $\gamma(x)\in\{0,1\}$.\end{itemize}
If $\alpha$ and $\beta$ are two partial assignments to variables from the set $X$, we say that a generalized partial assignment $\alpha\joint\beta:X\to\{0,1,*,?\}$ is a joint assignment iff:
\begin{itemize}\item if $\alpha(x)=a$ and $\beta(x)\in\{a,*\}$, then $(\alpha\joint\beta)(x)=a$;
\item if $\beta(x)=a$ and $\alpha(x)\in\{a,*\}$, then $(\alpha\joint\beta)(x)=a$;
\item if $\alpha(x)=a$ and $\beta(x)=1-a$, then $(\alpha\joint\beta)(x)=?$;
\item if $\alpha(x)=\beta(x)=*$, then $(\alpha\joint\beta)(x)=*$,
\end{itemize}
where $a\in\{0,1\}$.
\begin{result}{Theorem 4 (Max-Flow Min-Cut [8])}
Let $G:=(V,E)$. For any $s,t\in V$ the maximum value of an $s$-$t$ flow is equal to the minimum capacity over all $s$-$t$ cuts.
\end{result}
Let $X:=\{x_1,\ldots,x_n\}$ be a set of propositional variables and $O$ be a finite set. A branching program is a directed acyclic graph with one source. Every vertex of the graph is labeled by a variable from $X$, or by an element of the set $O$ with respect to the following properties:
\begin{itemize}\item if a vertex is labeled by $o\in O$, then it is a sink;
\item if a vertex is labeled by a variable, then it has exactly two outgoing edges: one edge is labeled by $0$ and the other one is labeled by $1$.\end{itemize}
Every branching program $B$ defines a function $f_B:\{0,1\}^n\to O$. We assume that every input $z\in\{0,1\}^n$ induces a path from source to sink in a natural way. If this path ends in a vertex with a label $o\in O$ then we define $f_B(z):=o$.
We say that $B$ is a branching program for the relation $S\subseteq\{0,1\}\times O$ iff $f_B$ is consistent with $S$: namely if $f_B(z)=o$ then $(z,o)\in S$.
Let $D$ be a branching program and $v$ is a node in it. The subprogram of $D$ with the root $v$ we denote by $D(v)$ and define as a subgraph of $D$ that is reachable from $v$. Also for a partial assignment $\rho$ we define a branching program $D|_\rho$ as the following transformation applied to $D$:
\begin{itemize}\item for each variable $y$ to which $\rho$ assigns a value $a$, contract edges $y=a$ and delete edges $y=\neg a$;
\item delete all vertices that are unreachable from the root.\end{itemize}
These operations only decrease the size of the program.
If $p$ is a consistent path in a branching program, we denote a partial assignment that corresponds to this path by $\tau_p$.
\begin{result}{Definition 5}
Let $B$ be a branching program. We say that $B$ is a (syntactic) read-once branching program or $1$-BP iff on every path from the source to a sink we can see each variable at most once.
We say that $B$ is a $(1,+k)$-BP iff on every path $p$ from the source to a sink there is a set of variables $X_p$ of size at most $k$ such that all other variables appear in $p$ at most once.
And we can twist this definition a little bit and say that $B$ is a semantic $(1,+k)$-BP iff on every consistent path from $p$ source to sink there is a set of variables $X_p$ of size at most $k$ such that all other variables appear in $p$ at most once.
\end{result}
\begin{result}{Lemma 8}
Let $D$ be a $(1,+k)$-BP for $\Search_\varphi$, $p$ be a consistent path from the root to some node $v$. If $p$ has a variable $x$ queried more than one time on it then $D(v)|_{\tau_p}$ is a $(1,+(k-1))$-BP for the $\Search_{\varphi|_{\tau_p}}$. The result holds for both: semantic and syntactic models.
\end{result}
\begin{proof}
A program $D(v)|_{\tau_p}$ is a program for the $\Search_{\varphi|_{\tau_p}}$ by the correctness of the program $D$. Consider a path $s$ in $D$ from $v$ to some leaf. Let $X_s$ be a set of variables that are queried more than one time on $s$. If $|X_s|=k$ and $x\notin X_s$, the path $ps$ has at least $k+1$ variables that are queried more that one time. This is a contradiction. If $|X_s|=k$ and $x\in X_s$, note that in $D(v)|_{\tau_p}$ we contract all edges that correspond to the $x$ variable and hence we transform this path into a path with at most $k-1$ repetitions.
\end{proof}
\section*{Original Section 3: Expanders (pp.17:6--7)}
We are given a graph $G:=(V,E)$. For two subsets of vertices $A,B$ we write $E(A,B)$ to denote the set of pairs $(v,e)$ where $v\in A$, $e$ is an edge that is incident to $v$ and $e$ connects $v$ with some vertex in $B$. We will think about it as about set of edges between $A$ and $B$, but if $A$ and $B$ intersect we count edges within intersection twice. We also use a shortcut notations $E(S):=E(S,V)$ and $\overline S:=V\setminus S$. If the graph we consider is unclear from the context we specify it as a subscript: $E_G(A,B)$.
\begin{result}{Remark 9}
Assuming that $G$ is $\Delta$-regular graph this definition allows us to use natural equalities:
\[|E(S)|=\Delta|S|;\qquad |E(A,\overline A)|=\Delta|A|-|E(A,A)|.\]
\end{result}
We write $N_G(v)$ to denote the set of neighbours of $v$ in the graph $G$. We extend this notion to sets and denote by $N_G(S):=\{v\mid\exists u\in S,(u,v)\in E\}$ the neighbourhood of a set of vertices $S\subseteq V$.
A graph $G:=(V,E)$ is an $(n,\Delta,\alpha)$-algebraic expander (or just expander), if:
\begin{itemize}\item $|V|=n$;
\item the degree of any vertex $v\in V$ equals $\Delta$;
\item the absolute value of the second largest eigenvalue of the adjacency matrix of $G$ is at most $\alpha\Delta$.\end{itemize}
\begin{result}{Lemma 10 (Mixing Lemma [4])}
Let $G:=(V,E)$ be an $(n,\Delta,\alpha)$-expander. For any two subsets $A,B\subseteq V$ the following holds:
\[\left||E(A,B)|-\frac{\Delta|A||B|}{n}\right|\le\alpha\Delta\sqrt{|A||B|}.\]
\end{result}
We also need combinatorial edge expansion. We say that $G:=(V,E)$ satisfies $(r,\beta)$-(edge) expansion property for some $r,\beta>0$, if for all $S\subseteq V$ of size at most $r$ holds $E(S,\overline S)\ge\beta\Delta|S|$. The Mixing Lemma says that any expander graph satisfies expansion property for suitable parameters.
\begin{result}{Corollary 11}
If $G:=(V,E)$ is an $(n,\Delta,\alpha)$-expander, then for any $0<\beta<1-\alpha$ the graph $G$ satisfy $((1-\alpha-\beta)n,\beta)$-expansion property.
\end{result}
\begin{proof}
Consider some $A\subseteq V$ of size at most $(1-\alpha-\beta)n$. Note that $|E(A,\overline A)|=\Delta|A|-|E(A,A)|$. By Mixing Lemma:
\[|E(A,A)|\le\frac{\Delta|A|^2}{n}+\alpha\Delta|A|=\Delta|A|\left(\frac{|A|}{n}+\alpha\right)\le\Delta|A|(1-\beta).\]
Hence $|E(A,\overline A)|\ge\beta\Delta|A|$ by Remark 9.
\end{proof}
The ``vertex analog'' of the next proposition is well known in the literature (for example [9]). We turn it into edge version.
\begin{result}{Proposition 12}
Let $G:=(V,E)$ be a graph of degree $\Delta$. If $G$ satisfies $(r,\beta)$-expansion property then for any set $S\subseteq V$ of size $k\le r$ there is an enumeration $v_1,v_2,\ldots,v_k\in S$ and a sequence $R_1,\ldots,R_k\subseteq E(S)$ such that:
\[R_i=E(\{v_i\},V\setminus\{v_1,v_2,\ldots,v_i\});\qquad |R_i|\ge\beta\Delta.\]
\end{result}
\begin{proof}
We create this sequence in reversed order. Since $|S|\le r$, it holds that $|E(S,\overline S)|\ge\beta\Delta|S|$ and there is a vertex $v_k\in S$ such that $|E(\{v_k\},\overline S)|\ge\beta\Delta$. Let $R_k:=|E(\{v_k\},\overline S)|$, and repeat the process for $S\setminus\{v_k\}$.
\end{proof}
